% 第3章第5节：本章小结。

\section{本章小结}
\label{sec:ose-summary}

章首的积分变换$Tf=\int(1+xz)f\,d\mu$只记住输入的两个矩，$\ell^2$上的$A(x_n)=(x_n/n)$则不丢掉任何方向，只是越往后压得越狠。本章围绕这两个例子，从三个角度描述一个变换：它能否用有限个方向统一逼近，各个方向被保留或压缩多少，以及从输出能否稳定地恢复输入。

紧性回答第一个问题。紧算子把有界输入族压成可以抽出收敛子列的输出族，并把弱收敛的输入变成范数收敛的输出。紧算子的算子范数极限仍是紧算子；在Hilbert空间中，紧算子恰好是有限秩算子的算子范数极限。这里的关键是“统一”：每个函数都能被截断逼近，并不意味着同一个有限维空间能以给定精度覆盖整个单位球的像，单位算子就是反例。

谱与奇异值回答第二个问题。无限维中谱不只由特征值组成，$A$的谱点$0$不对应任何被消去的方向，而对应一列越来越扁的方向和不闭的像。紧自伴算子有至多可数个实的非零特征值，只能在零处聚集，特征方向与核一起张满整个空间；一般紧算子用输入、输出两组方向写成奇异值分解，截断前$m$项的误差恰为$\sigma_{m+1}$，且是所有秩不超过$m$的逼近中最小的。这些展开都按Hilbert空间的范数成立，不自动给出逐点收敛。

积分核把上述结论具体化。核与测度一起确定算子；平方可积核给出有界的紧算子，证明的桥梁是用矩形示性函数的有限和逼近核；交换变量并取共轭得到伴随核，对称核给出自伴算子。但对称不保证正性，平方可积核作为等价类也谈不上点值和有限点集上的Gram矩阵。

第三个问题在两类方程上有不同的答案。第二类方程$(I-\lambda T)f=g$保留了单位算子，小范数时Neumann级数和压缩迭代直接给出解与误差界；对紧算子，Fredholm二择保证像闭合，于是要么唯一可解且稳定，要么在正交相容条件下可解、解差一个有限维核。第一类方程$Tf=g$还要满足Picard条件，而小奇异值使最小范数解对数据误差无界地敏感。几种判断不能互相替代：正交条件只在像闭合时判断精确可解，可解不等于某个迭代收敛，有解也不等于稳定。

这些认识是后续几类问题的起点。第~\ref{chap:rkhs}章给积分核加上连续性与半正定性，研究点值读取连续的函数空间，并在那里建立Mercer展开的逐点与一致收敛；带噪数据下怎样舍去或压低小奇异值方向，是优化中正则化的问题，具体的微分方程演化与信号重建则分别属于卷三的动力系统与信号部分。本章只保留算子方向与积分方程的共同基础，不预先引入这些任务各自的目标和算法。
